\begin{table}[t]
    \centering
    \resizebox{\textwidth}{!}{%
    \begin{tabular}{@{}lcccl@{}}
        \toprule
        \textbf{Ablation (vs.\ \texttt{D1\_main})} & \textbf{direction} & \textbf{effect (Cohen's $d_z$)} & \textbf{$p$} & \textbf{reading} \\
        \midrule
        single-stage vs.\ two-stage        & $\downarrow\downarrow\downarrow$ converges & $-2.56$ & $5\times10^{-4}$ & \textbf{the} decisive factor \\
        temperature $1.0$ vs.\ $0$         & $\uparrow$ more churn / chaotic & $+2.09$ & $5\times10^{-4}$ & default temp destroys even relaxation ($R^2{=}0.27$) \\
        cross-model reviewer vs.\ self     & $\downarrow$ less churn & $-1.17$ & $5\times10^{-3}$ & a different reviewer eases pressure, still no FP \\
        lenient vs.\ structured framing    & $\downarrow$ less churn & $-1.14$ & $7\times10^{-3}$ & less to fix $\Rightarrow$ smaller edits, still no FP \\
        harsh vs.\ structured framing      & $\uparrow$ (n.s.) & $+0.45$ & $0.13$ & trend toward more churn \\
        unconstrained vs.\ length-preserving & $\uparrow$ (n.s.) & $+0.19$ & $0.62$ & length guidance barely helps \\
        \bottomrule
    \end{tabular}
    }
    \caption{Determinants of (non-)convergence: paired comparisons against \texttt{D1\_main}
    (Wilcoxon on final $\Delta$). No manipulation produces convergence; the single- vs.\
    two-stage contrast dominates every within-two-stage factor. Negative $d_z$ means less
    churn than \texttt{D1\_main}. n.s.\ = not significant.}
    \label{tab:ablations}
\end{table}
